% Setting up the document and importing necessary packages
\documentclass[12pt,a4paper]{article}

% Chinese support packages
\usepackage{xeCJK}
\usepackage{fontspec}
\setCJKmainfont{Noto Serif CJK TC}
\setCJKsansfont{Noto Sans CJK TC}
\setCJKmonofont{Noto Sans CJK TC}

% Other necessary packages
\usepackage{geometry}
\usepackage{titlesec}
\usepackage{enumitem}
\usepackage{color}
\usepackage{colortbl}
\usepackage{tcolorbox}
\usepackage{booktabs}
\usepackage{longtable}
\usepackage{array}
\usepackage{graphicx}
\usepackage{tikz}
\usepackage{mdframed}
\usepackage{fancyhdr}
\usepackage{hyperref}
\usepackage{multicol}
\usepackage{datetime2}
\usepackage{natbib} % For bibliography

\hypersetup{
    colorlinks=true,
    linkcolor=nephro_blue,
    citecolor=nephro_blue,
    urlcolor=nephro_blue,
    pdfborder={0 0 0}
}

% Page setup
\geometry{
    top=2.5cm,
    bottom=2.5cm,
    left=2cm,
    right=2cm,
    headheight=15pt
}

% Color definitions (Nephrology themed colors)
\definecolor{nephro_blue}{RGB}{0,102,204}
\definecolor{nephro_cyan}{RGB}{0,153,204}
\definecolor{nephro_gray}{RGB}{120,120,120}
\definecolor{highlight_yellow}{RGB}{255,250,205}

% Title format settings
\titleformat{\section}[block]{\Large\bfseries\color{nephro_blue}}{\thesection}{10pt}{}
\titleformat{\subsection}[block]{\large\bfseries\color{nephro_cyan}}{\thesubsection}{8pt}{}
\titleformat{\subsubsection}[block]{\normalsize\bfseries\color{nephro_gray}}{\thesubsubsection}{6pt}{}

% Custom environment
\newmdenv[
    linecolor=nephro_blue,
    backgroundcolor=highlight_yellow,
    linewidth=2pt,
    topline=true,
    bottomline=true,
    leftline=true,
    rightline=true,
    innertopmargin=10pt,
    innerbottommargin=10pt,
    innerleftmargin=10pt,
    innerrightmargin=10pt
]{importantbox}

\newcommand{\note}[1]{
    \par
    \noindent
    \begin{tcolorbox}[
        colback=nephro_cyan!5!white,
        colframe=nephro_cyan,
        title=\textbf{重點提醒},
        fonttitle=\bfseries,
        boxrule=0.5pt,
        arc=3pt,
        left=6pt,
        right=6pt,
        top=6pt,
        bottom=6pt
    ]
    #1
    \end{tcolorbox}
}

% Header and footer settings
\pagestyle{fancy}
\fancyhf{}
\fancyhead[L]{\textcolor{nephro_blue}{內科部住院醫師教學文件}}
\fancyhead[R]{\textcolor{nephro_cyan}{腎臟內科EPAs}}
\fancyfoot[C]{\textcolor{nephro_gray}{\thepage}}
\renewcommand{\headrulewidth}{1pt}
\renewcommand{\headrule}{\hbox to\headwidth{\color{nephro_blue}\leaders\hrule height \headrulewidth\hfill}}

% Date format
\DTMsetstyle{yyyy-mm-dd}
\newcommand{\mydate}{\DTMtoday}

% Document start
\begin{document}

% Title page with table of contents
\begin{titlepage}
    \centering
    {\LARGE\textcolor{nephro_blue}{\textbf{內科部住院醫師教學文件}}\par}
    \vspace{1cm}
    {\huge\bfseries 腎臟內科可信賴專業活動EPAs\par}
    \vspace{0.5cm}
    {\Large\textcolor{nephro_cyan}{\textit{Nephrology Entrustable Professional Activities}}\par}
    \vspace{1cm}
    
    \begin{tcolorbox}[
        colback=nephro_cyan!5!white,
        colframe=nephro_cyan,
        width=0.8\textwidth,
        arc=10pt,
        boxrule=1pt
    ]
    \centering
    {\large\textbf{目標對象}\par}
    \vspace{0.3cm}
    內科部住院醫師R1-R3\par
    腎臟內科專科訓練醫師\par
    \vspace{0.5cm}
    {\large\textbf{文件版本}\par}
    \vspace{0.3cm}
    第1.0版(10個EPAs)\par
    發布日期:\mydate\par
    \end{tcolorbox}

    \vfill
    
    {\large 腎臟內科 陳長江醫師\\
    適用於CCC評量系統\par}
    
    \vspace{0.5cm}
\end{titlepage}

\newpage
\tableofcontents
\newpage

% Main content
\section{腎臟內科EPAs概述}

\subsection{EPA定義}
Entrustable Professional Activities (EPA) 是指可信賴委託給住院醫師執行的專業活動單元。腎臟內科EPAs涵蓋了腎臟疾病診斷、治療與照護的核心能力,是所有腎臟內科醫師必須精熟的專業技能。

\vspace{0.5cm}
\note{
腎臟內科EPAs評量著重於能否安全、穩定地執行完整的腎臟疾病診療流程,並能整合資訊做出適當臨床決策。每個EPA都包含明確的里程碑發展階段,從初學者到專家的完整能力建構。本版本包含10個核心EPA,從基礎的病史詢問與身體檢查,到複雜的腎臟替代療法與跨科部協調,提供完整的腎臟疾病診療能力架構。
}

\subsection{學習目標}
完成腎臟內科EPAs後,住院醫師應能夠:
\begin{itemize}[nosep]
    \item 熟練進行腎臟疾病病史詢問與風險評估。
    \item 準確執行腎臟泌尿系統身體檢查。
    \item 熟悉腎臟超音波基礎操作與判讀。
    \item 獨立判讀酸鹼及電解質檢驗報告。
    \item 診斷與處理急性腎損傷。
    \item 診斷與管理慢性腎臟疾病。
    \item 處置體液及電解質不平衡。
    \item 認識並評估腎臟替代療法。
    \item 評估腎臟切片檢查適應症及風險。
    \item 協調末期腎病病人跨科部整合照護。
\end{itemize}

\subsection{委託等級}
\begin{table}[h]
    \centering
    \caption{EPA委託等級定義}
    \begin{tabular}{p{2cm} p{3cm} p{9cm}}
        \toprule
        \textbf{等級} & \textbf{委託程度} & \textbf{描述} \\
        \midrule
        Level 1 & 觀察學習 & 僅能觀察他人執行,無法獨立操作 \\
        Level 2 & 直接監督 & 可在主治醫師直接監督下執行 \\
        Level 3 & 間接監督 & 可獨立執行,但需回報並接受指導 \\
        Level 4 & 監督可得 & 可獨立執行,監督者隨時可得 \\
        Level 5 & 完全獨立 & 可獨立執行並指導他人 \\
        \bottomrule
    \end{tabular}
\end{table}

\newpage
\section{腎臟內科10大EPAs詳述}

\begin{longtable}{|p{1.2cm}|p{3.8cm}|p{8.5cm}|}
\caption{腎臟內科EPAs完整架構} \\
\hline
\rowcolor{nephro_cyan!20}
\textbf{EPA} & \textbf{專業活動名稱} & \textbf{核心能力要素與里程碑} \\
\hline
\endfirsthead

\multicolumn{3}{c}%
{{\bfseries \tablename\ \thetable{} -- 接續上頁}} \\
\hline
\rowcolor{nephro_cyan!20}
\textbf{EPA} & \textbf{專業活動名稱} & \textbf{核心能力要素與里程碑} \\
\hline
\endhead

\hline \multicolumn{3}{|r|}{{接續下頁}} \\ \hline
\endfoot

\hline \hline
\endlastfoot

\rowcolor{nephro_blue!10}
\multicolumn{3}{|c|}{\textbf{EPA 1: 腎臟疾病病史詢問與風險評估}} \\
\hline
\textbf{描述} & \multicolumn{2}{|p{12.5cm}|}{能夠獨立進行腎臟疾病相關的完整病史詢問,包括症狀評估、危險因子識別、用藥史調查,並能初步評估腎臟疾病風險。} \\
\hline
\textbf{核心能力} & 症狀評估 & 尿量變化、水腫、血尿、泡沫尿、腰痛的詳細詢問 \\
\cline{2-3}
& 危險因子識別 & 高血壓、糖尿病、自體免疫疾病、藥物史 \\
\cline{2-3}
& 功能狀態評估 & 尿毒症狀、體液平衡、電解質症狀評估 \\
\cline{2-3}
& 藥物史 & 腎毒性藥物使用史、過敏史、順從性評估 \\
\cline{2-3}
& 家族史 & 多囊性腎病、遺傳性腎病、腎結石家族史 \\
\hline
\textbf{Level 1} & Novice & 能詢問基本腎臟疾病症狀,需要監導者引導完成詢問 \\
\hline
\textbf{Level 2} & Advanced Beginner & 能系統性完成腎臟病史詢問,使用適當的症狀評估工具 \\
\hline
\textbf{Level 3} & Competent & 能獨立完成複雜病史詢問,整合症狀與危險因子進行初步風險分層 \\
\hline
\textbf{Level 4} & Proficient & 能在困難情況下獲取準確病史,整合多重共病資訊 \\
\hline
\textbf{Level 5} & Expert & 能指導他人進行病史詢問,識別罕見症狀表現 \\
\hline

\rowcolor{nephro_blue!10}
\multicolumn{3}{|c|}{\textbf{EPA 2: 腎臟泌尿系統身體檢查}} \\
\hline
\textbf{描述} & \multicolumn{2}{|p{12.5cm}|}{能夠熟練且準確地進行腎臟泌尿系統的完整身體檢查,包括腎臟觸診、膀胱檢查、水腫評估,並能正確解讀檢查發現。} \\
\hline
\textbf{核心能力} & 腎臟觸診 & 腎臟大小、壓痛、多囊腎觸診 \\
\cline{2-3}
& 水腫評估 & 下肢水腫、腹水、肺水腫評估 \\
\cline{2-3}
& 血壓測量 & 標準測量技術、姿勢性血壓變化 \\
\cline{2-3}
& 體液狀態 & 脫水徵象、體液過多徵象評估 \\
\cline{2-3}
& 特殊檢查 & 肋脊角壓痛、膀胱叩診 \\
\hline
\textbf{Level 1} & Novice & 能進行基本腎臟觸診,識別明顯水腫 \\
\hline
\textbf{Level 2} & Advanced Beginner & 能識別常見異常身體發現,進行系統性腎臟檢查 \\
\hline
\textbf{Level 3} & Competent & 能區分各種異常身體發現,準確評估體液狀態,整合多項檢查發現 \\
\hline
\textbf{Level 4} & Proficient & 能識別複雜身體發現特徵,檢查技術精準且有效率 \\
\hline
\textbf{Level 5} & Expert & 能偵測細微的異常發現,整合檢查發現與臨床表現 \\
\hline

\rowcolor{nephro_blue!10}
\multicolumn{3}{|c|}{\textbf{EPA 3: 腎臟超音波基礎操作與判讀}} \\
\hline
\textbf{描述} & \multicolumn{2}{|p{12.5cm}|}{能夠進行基礎腎臟超音波檢查,包括設備操作、解剖定位、基本測量與主要結構評估,並能識別重要的病理發現。} \\
\hline
\textbf{核心能力} & 檢查操作 & 超音波操作、病患指導、品質確保 \\
\cline{2-3}
& 基本測量 & 腎臟大小、皮質厚度、輸尿管擴張 \\
\cline{2-3}
& 結構評估 & 腎實質、集尿系統、血流評估 \\
\cline{2-3}
& 病變識別 & 囊腫、結石、水腎、實質病變 \\
\cline{2-3}
& 臨床關聯 & 超音波發現與疾病嚴重度關聯 \\
\hline
\textbf{Level 1} & Novice & 能協助進行基本超音波檢查,了解基本解剖結構 \\
\hline
\textbf{Level 2} & Advanced Beginner & 能獨立操作標準腎臟超音波檢查,進行基本判讀 \\
\hline
\textbf{Level 3} & Competent & 能完成標準腎臟超音波檢查,準確識別常見病變 \\
\hline
\textbf{Level 4} & Proficient & 能進行進階超音波技術,準確定量評估 \\
\hline
\textbf{Level 5} & Expert & 成為超音波技術專家,能教導複雜技術 \\
\hline

\rowcolor{nephro_blue!10}
\multicolumn{3}{|c|}{\textbf{EPA 4: 酸鹼及電解質檢驗報告判讀}} \\
\hline
\textbf{描述} & \multicolumn{2}{|p{12.5cm}|}{能夠獨立判讀酸鹼及電解質檢驗報告,包括動脈血液氣體分析、血清電解質,識別正常與異常變化,並提出臨床相關性分析。} \\
\hline
\textbf{核心能力} & 酸鹼判讀 & pH、HCO3、PaCO2解讀,代償機轉判斷 \\
\cline{2-3}
& 電解質異常 & 鈉、鉀、鈣、磷、鎂異常識別 \\
\cline{2-3}
& 病因分析 & 酸鹼及電解質異常病因推理 \\
\cline{2-3}
& 嚴重度評估 & 緊急處置需求評估 \\
\cline{2-3}
& 治療指引 & 矯正策略建議 \\
\hline
\textbf{Level 1} & Novice & 能識別明顯異常數值,知道基本正常範圍 \\
\hline
\textbf{Level 2} & Advanced Beginner & 能判讀常見酸鹼電解質異常,使用系統性判讀流程 \\
\hline
\textbf{Level 3} & Competent & 能獨立判讀複雜酸鹼電解質異常,整合臨床表現 \\
\hline
\textbf{Level 4} & Proficient & 能判讀複雜混合性酸鹼異常,準確推理病因 \\
\hline
\textbf{Level 5} & Expert & 成為酸鹼電解質判讀專家,能教導複雜判讀 \\
\hline

\rowcolor{nephro_blue!10}
\multicolumn{3}{|c|}{\textbf{EPA 5: 急性腎損傷診斷與處置}} \\
\hline
\textbf{描述} & \multicolumn{2}{|p{12.5cm}|}{能夠快速識別急性腎損傷,進行適當的診斷性檢查,評估嚴重度分級,並執行初步治療措施,包括體液管理與腎毒性藥物調整。} \\
\hline
\textbf{核心能力} & 快速識別 & AKI定義、KDIGO分期應用 \\
\cline{2-3}
& 診斷策略 & 腎前、腎性、腎後性AKI區分 \\
\cline{2-3}
& 嚴重度分級 & RIFLE、AKIN、KDIGO分級系統 \\
\cline{2-3}
& 體液管理 & 容積狀態評估、液體復甦策略 \\
\cline{2-3}
& 處置決策 & 腎毒性藥物停用、劑量調整、RRT評估 \\
\hline
\textbf{Level 1} & Novice & 能識別典型AKI症狀,知道基本檢查項目 \\
\hline
\textbf{Level 2} & Advanced Beginner & 能使用標準診斷流程,進行基本嚴重度評估 \\
\hline
\textbf{Level 3} & Competent & 能獨立診斷AKI,適當選擇治療策略 \\
\hline
\textbf{Level 4} & Proficient & 能處理複雜AKI情況,整合多重因素做決策 \\
\hline
\textbf{Level 5} & Expert & 成為AKI專家,能指導複雜個案處理 \\
\hline

\rowcolor{nephro_blue!10}
\multicolumn{3}{|c|}{\textbf{EPA 6: 慢性腎臟疾病診斷與管理}} \\
\hline
\textbf{描述} & \multicolumn{2}{|p{12.5cm}|}{能夠診斷慢性腎臟疾病,評估腎功能損傷程度,制定個別化治療計畫,包括藥物調整與非藥物治療,並監測疾病進展。} \\
\hline
\textbf{核心能力} & 診斷評估 & eGFR計算、蛋白尿評估、CKD分期 \\
\cline{2-3}
& 病因識別 & 糖尿病腎病變、高血壓腎病變等區分 \\
\cline{2-3}
& 功能評估 & 腎絲球濾過率、蛋白尿程度評估 \\
\cline{2-3}
& 藥物治療 & ACEI/ARB、降磷劑、貧血治療 \\
\cline{2-3}
& 整合照護 & 併發症預防、營養、疫苗接種建議 \\
\hline
\textbf{Level 1} & Novice & 能識別CKD症狀徵象,知道基本分期 \\
\hline
\textbf{Level 2} & Advanced Beginner & 能使用標準治療指引,進行基本藥物調整 \\
\hline
\textbf{Level 3} & Competent & 能獨立管理穩定CKD,調整藥物治療 \\
\hline
\textbf{Level 4} & Proficient & 能處理複雜CKD,評估進階治療選項 \\
\hline
\textbf{Level 5} & Expert & 成為CKD專家,發展個人化治療策略 \\
\hline

\rowcolor{nephro_blue!10}
\multicolumn{3}{|c|}{\textbf{EPA 7: 體液及電解質不平衡處置}} \\
\hline
\textbf{描述} & \multicolumn{2}{|p{12.5cm}|}{能夠診斷各種體液及電解質不平衡,評估嚴重度,制定矯正策略,監測治療效果與副作用,並提供適當的體液電解質管理。} \\
\hline
\textbf{核心能力} & 臨床評估 & 體液過多/不足、電解質異常症狀 \\
\cline{2-3}
& 病因分析 & 低血鈉、高血鉀等病因推理 \\
\cline{2-3}
& 嚴重度評估 & 緊急處置需求判斷 \\
\cline{2-3}
& 矯正策略 & 輸液種類選擇、矯正速率計算 \\
\cline{2-3}
& 監測管理 & 治療反應評估、併發症預防 \\
\hline
\textbf{Level 1} & Novice & 能識別明顯體液電解質異常,知道基本處理原則 \\
\hline
\textbf{Level 2} & Advanced Beginner & 能處理常見體液電解質異常,使用標準矯正公式 \\
\hline
\textbf{Level 3} & Competent & 能獨立處理大部分體液電解質異常,制定矯正計畫 \\
\hline
\textbf{Level 4} & Proficient & 能處理複雜混合性異常,整合多重因素 \\
\hline
\textbf{Level 5} & Expert & 成為體液電解質專家,發展個人化處置策略 \\
\hline

\rowcolor{nephro_blue!10}
\multicolumn{3}{|c|}{\textbf{EPA 8: 腎臟替代療法認識與評估}} \\
\hline
\textbf{描述} & \multicolumn{2}{|p{12.5cm}|}{能夠認識各種腎臟替代療法,包括血液透析、腹膜透析、腎臟移植,評估適應症,提供病患諮詢,並協助治療模式選擇。} \\
\hline
\textbf{核心能力} & 治療模式 & 血液透析、腹膜透析、腎臟移植認識 \\
\cline{2-3}
& 適應症評估 & RRT啟動時機、緊急透析適應症 \\
\cline{2-3}
& 血管通路 & 動靜脈瘻管、導管、腹膜透析導管 \\
\cline{2-3}
& 病患諮詢 & 治療選項說明、預期結果討論 \\
\cline{2-3}
& 併發症管理 & 透析相關併發症識別與處理 \\
\hline
\textbf{Level 1} & Novice & 能了解RRT基本概念,識別透析適應症 \\
\hline
\textbf{Level 2} & Advanced Beginner & 能評估RRT需求,提供基本病患諮詢 \\
\hline
\textbf{Level 3} & Competent & 能獨立評估RRT適應症,協助模式選擇 \\
\hline
\textbf{Level 4} & Proficient & 能處理複雜RRT評估,整合多重考量因素 \\
\hline
\textbf{Level 5} & Expert & 成為RRT專家,提供個人化治療建議 \\
\hline

\rowcolor{nephro_blue!10}
\multicolumn{3}{|c|}{\textbf{EPA 9: 腎臟切片檢查適應症及風險評估}} \\
\hline
\textbf{描述} & \multicolumn{2}{|p{12.5cm}|}{能夠評估腎臟切片檢查適應症,識別禁忌症,進行風險評估,提供病患諮詢,並了解檢查前後的管理原則。} \\
\hline
\textbf{核心能力} & 適應症評估 & 診斷性vs治療性切片決策 \\
\cline{2-3}
& 禁忌症識別 & 絕對與相對禁忌症判斷 \\
\cline{2-3}
& 風險評估 & 出血、感染、其他併發症風險 \\
\cline{2-3}
& 病患諮詢 & 檢查必要性、風險說明、同意書取得 \\
\cline{2-3}
& 術後管理 & 併發症監測、病理報告解讀 \\
\hline
\textbf{Level 1} & Novice & 能了解腎切片基本適應症,識別常見禁忌症 \\
\hline
\textbf{Level 2} & Advanced Beginner & 能評估基本適應症,進行初步風險評估 \\
\hline
\textbf{Level 3} & Competent & 能獨立評估腎切片適應症,完整風險評估 \\
\hline
\textbf{Level 4} & Proficient & 能處理複雜評估情況,整合多重風險因素 \\
\hline
\textbf{Level 5} & Expert & 成為腎切片評估專家,處理高風險個案 \\
\hline

\rowcolor{nephro_blue!10}
\multicolumn{3}{|c|}{\textbf{EPA 10: 末期腎病病人跨科部整合照護}} \\
\hline
\textbf{描述} & \multicolumn{2}{|p{12.5cm}|}{能夠協調末期腎病病人的整體照護,包括跨科部會診、多專科團隊合作、透析安排,並能有效溝通與教育病患及家屬。} \\
\hline
\textbf{核心能力} & 會診協調 & 心臟科、內分泌科、外科會診安排 \\
\cline{2-3}
& 團隊合作 & 多專科團隊會議參與、治療計畫整合 \\
\cline{2-3}
& 透析安排 & RRT模式選擇、血管通路安排 \\
\cline{2-3}
& 病患教育 & 疾病衛教、生活型態調整指導 \\
\cline{2-3}
& 品質改善 & 照護品質監測、改善方案執行 \\
\cline{2-3}
& 家屬溝通 & 病情解釋、治療選項討論 \\
\cline{2-3}
& 倫理考量 & 醫療倫理、安寧照護決策 \\
\hline
\textbf{Level 1} & Novice & 能參與基本病患照護,了解團隊角色 \\
\hline
\textbf{Level 2} & Advanced Beginner & 能執行基本會診,參與團隊討論 \\
\hline
\textbf{Level 3} & Competent & 能獨立協調病患照護,制定整合計畫 \\
\hline
\textbf{Level 4} & Proficient & 能整合複雜多科照護,領導團隊討論 \\
\hline
\textbf{Level 5} & Expert & 成為照護協調專家,改善照護品質 \\
\hline

\end{longtable}

\newpage
\section{評量方法與標準}

\subsection{CCC評量架構}
本腎臟內科EPAs採用Case-based Clinical Competency (CCC)評量方式:
\begin{itemize}[nosep]
    \item 真實臨床情境評量
    \item 多次觀察與回饋
    \item 漸進式委託決策
    \item 360度回饋機制
\end{itemize}

\begin{importantbox}
\textbf{特別提醒:} 腎臟內科EPAs的評量重點在於漸進式能力發展,而非單次完美表現。鼓勵從錯誤中學習,持續改進。每個EPA都需要在真實臨床情境中反復練習與評量。
\end{importantbox}

\subsection{評量工具對應表}
\begin{table}[h]
    \centering
    \caption{各EPA推薦評量工具}
    \begin{tabular}{p{4.5cm} p{4cm} p{5.5cm}}
        \toprule
        \textbf{EPA類型} & \textbf{主要評量工具} & \textbf{輔助評量方法} \\
        \midrule
        技能操作類 & DOPS & 直接觀察、技能檢核表 \\
        (EPA 3, 9) & (Direct Observation of Procedural Skills) & 同儕評量、自我評量 \\
        \midrule
        臨床推理類 & Mini-CEX & 病例討論、口試 \\
        (EPA 1, 4, 5, 6, 7, 8) & (Mini-Clinical Evaluation Exercise) & 病歷撰寫評量、模擬情境 \\
        \midrule
        身體檢查類 & DOPS & 標準化病人、技能測驗 \\
        (EPA 2) & (Direct Observation of Procedural Skills) & 影片回顧、同儕觀察 \\
        \midrule
        整合照護類 & Multi-source Feedback & 跨科團隊評量 \\
        (EPA 10) & (360-degree Assessment) & 病人滿意度調查 \\
        \bottomrule
    \end{tabular}
\end{table}

\subsection{評量頻率建議}
\begin{itemize}[nosep]
    \item \textbf{每月評量}:選擇3-4個EPA重點評量
    \item \textbf{季度檢討}:全面檢視所有10個EPA進展
    \item \textbf{年度認證}:EPA達標認證與進階規劃
    \item \textbf{即時回饋}:發現學習需求時立即介入
\end{itemize}

\vspace{0.5cm}
\note{
腎臟內科EPAs的學習是一個持續的過程。建議住院醫師建立個人學習檔案,記錄每次評量結果與改善計畫。主治醫師應提供及時具體的回饋,協助學習者達到下一個里程碑。
}

\section{實施建議與學習資源}

\subsection{月度晨會Mini-OSCE整合建議}
\begin{table}[h]
    \centering
    \caption{月度EPA評量主題安排}
    \begin{tabular}{p{2.5cm} p{4.5cm} p{7cm}}
        \toprule
        \textbf{月份} & \textbf{重點EPAs} & \textbf{評量重點} \\
        \midrule
        第1個月 & EPA 1, 2 & 基礎評估技能建立 \\
        第2個月 & EPA 3, 4 & 影像與檢驗判讀能力 \\
        第3個月 & EPA 5, 6 & 急慢性腎病處理 \\
        第4個月 & EPA 7 & 體液電解質管理 \\
        第5個月 & EPA 8, 9 & 進階治療評估 \\
        第6個月 & EPA 10 & 整合照護與協調能力 \\
        第7-12個月 & 綜合評量與進階訓練 & 複雜個案整合處理能力 \\
        \bottomrule
    \end{tabular}
\end{table}

\subsection{推薦學習資源}
學員可參考以下文獻深入學習:

\textbf{基礎教科書:}
\begin{itemize}
    \item Brenner \& Rector's The Kidney. 11th edition.
    \item Oxford Textbook of Clinical Nephrology. 4th edition.
    \item Comprehensive Clinical Nephrology. 6th edition.
\end{itemize}

\textbf{臨床指引:}
\begin{itemize}
    \item KDIGO Clinical Practice Guidelines for Acute Kidney Injury
    \item KDIGO Clinical Practice Guidelines for CKD Evaluation and Management
    \item KDOQI Clinical Practice Guidelines for Nutrition in CKD
    \item KDIGO Clinical Practice Guidelines for Anemia in CKD
    \item KDIGO Clinical Practice Guidelines for CKD-MBD
\end{itemize}

\subsection{自我評估工具}
\begin{itemize}[nosep]
    \item 定期記錄學習歷程與反思
    \item 尋求多方回饋(主治醫師、護理師、病人)
    \item 參與同儕學習與討論
    \item 利用模擬教學工具練習
    \item 建立個人學習檔案
    \item 參與多專科團隊會議
    \item 定期檢視腎臟超音波技能
\end{itemize}

\section{總結}

腎臟內科EPAs涵蓋了腎臟疾病診療的完整核心能力,從基礎的病史詢問與身體檢查,到複雜的腎臟替代療法評估、腎切片風險評估,以及整合性的跨科部協調照護。每個EPA都有清楚的能力描述與里程碑發展路徑,協助住院醫師循序漸進地建立專業能力。

\begin{importantbox}
\textbf{關鍵訊息:} 腎臟內科EPAs的核心在於「可委託性」與「完整性」。這10個EPA代表腎臟內科醫師必須具備的基礎診療能力。當您能夠安全、獨立地執行各項腎臟疾病診療活動,並做出合理的臨床決策時,就達到了相對應EPA的目標。記住,每一次的臨床實務都是學習的機會,透過持續的實作、反思與改進,才能成為一位優秀的腎臟內科醫師。
\end{importantbox}

% References section
\section{參考文獻}
\begin{thebibliography}{10}

\bibitem{brenner2019}
Yu, A. S. L., et al. (2019). \textit{Brenner \& Rector's The Kidney}. 11th edition. Elsevier.

\bibitem{oxford2015}
Turner, N., et al. (2015). \textit{Oxford Textbook of Clinical Nephrology}. 4th edition. Oxford University Press.

\bibitem{olle2013}
ten Cate, O. (2013). Nuts and bolts of entrustable professional activities. \textit{Journal of Graduate Medical Education}, 5(1), 157-158.

\bibitem{kdigo2012aki}
KDIGO AKI Work Group. (2012). KDIGO Clinical Practice Guideline for Acute Kidney Injury. \textit{Kidney International Supplements}, 2(1), 1-138.

\bibitem{kdigo2024ckd}
KDIGO 2024 Clinical Practice Guideline for the Evaluation and Management of Chronic Kidney Disease. \textit{Kidney International}, 105(4S), S117-S314.

\bibitem{kdoqi2020nutrition}
Ikizler, T. A., et al. (2020). KDOQI Clinical Practice Guideline for Nutrition in CKD: 2020 Update. \textit{American Journal of Kidney Diseases}, 76(3 Suppl 1), S1-S107.

\bibitem{johnson2018}
Johnson, R. J., Floege, J., \& Feehally, J. (2018). \textit{Comprehensive Clinical Nephrology}. 6th edition. Elsevier.

\end{thebibliography}

\end{document}